\section{A spectral memory law for two quantum detectors}\label{sec:two-detectors}
We apply the finite-memory framework without a balancing group action or reset to a prescribed reference state. The measurements disturb the state: this is first detection, not estimation of an untouched pre-command parameter. Quantum first-detection laws, dark states and probing-time optimization have substantial antecedents \cite{Kessler21,DelVecchio26}. Our question is the exact paid-call versus classical-label frontier when the observer chooses between two specified detectors.

Let $P,Q$ be orthogonal projections on a finite-dimensional Hilbert space. Detector $A$ has failure map $\rho\mapsto P\rho P$ and success map $\rho\mapsto(I-P)\rho(I-P)$; detector $B$ uses $Q$ in the same way. Every call costs one. There is no intervening evolution, additional measurement, reset or coherent controller memory. A failure does not permit commitment. The first success requires immediate productive commitment, with one constant terminal decoder value. The initial density matrix is arbitrary and unknown. An autonomous controller has a prescribed initial live label, fixed command rows and post-failure transitions, with no free clock or persistent external seed. Write $V_j$ for the least worst-initial-state expected call count using at most $j$ live labels, and $V_*$ for the infimum over all visible-history policies.

\begin{theorem}[Exact two-detector resource law]\label{thm:detector-memory}
Suppose $P,Q\ne0$ and $c=\|PQ\|<1$. In every finite dimension,
\begin{equation}\label{eq:detector-global}
 V_*=1+\frac1{2(1-c)}.
\end{equation}
It is attained by choosing the first detector fairly and then switching after every failure. The one-live-label value is
\begin{equation}\label{eq:detector-one}
 V_1=1+\frac1{1-c}=2V_*-1,
\end{equation}
and a fair independent detector choice at every call attains it. With a fixed initial label, three live labels are necessary and sufficient to attain $V_*$. The two-live-label optimum is attained and satisfies
\begin{equation}\label{eq:detector-two}
 V_*<V_2\le1+\frac1{1-c^2};
\end{equation}
its exact value is determined in Theorem~\ref{thm:two-label-exact} below.

If an explicitly randomized initial distribution over counted labels is instead part of the interface, two live labels are necessary and sufficient to attain $V_*$. Under event stopping the stated live counts are also total-label counts, since the constant terminal value may reuse a live label. State stopping requires one additional disjoint terminal label.

If either failure projection is zero, one call to that detector gives value one. If $c=1$, their ranges have a common nonzero dark subspace and no policy is proper for every initial state. The robust cost is then infinite.
\end{theorem}
\begin{proof}
We give the required spectral reduction explicitly. Diagonalize $PQP$ on $\operatorname{ran}P$. For an eigenvector $u$ with eigenvalue $q\in(0,1)$, put $v=Qu/\sqrt q$. Then $\|v\|=1$ and $\langle u,v\rangle=\sqrt q$. Their span is invariant under both projections, which restrict to the rank-one projections onto $u,v$. Distinct orthonormal eigenvectors give orthogonal such two-dimensional blocks, including after resolving multiplicities. The zero eigenspace gives $\operatorname{ran}P\cap\ker Q$; the remaining part of $\operatorname{ran}Q$ is orthogonal to $P$. The orthogonal complement of the two ranges lies in both kernels. Eigenvalue one is exactly a common failure range. This is the classical two-subspace geometry \cite{Halmos69}, proved here in the form used below. The largest nontrivial squared overlap is $q=c^2$ when $c>0$.

After a failure of $A$, a second immediate $A$ fails with probability one and leaves the state unchanged. The analogous statement holds for $B$. Such a repetition gives no information. In a full-history policy it can be deleted, with its private coins and deterministic failure simulated internally, leaving all subsequent physical states and detector choices coupled and weakly reducing calls. Infinite repetitions give infinite cost and cannot improve the optimum. Thus, after its first choice, a pointwise optimal policy switches after every failure.

On a block of squared overlap $q$, the post-failure pure state switches between $u$ and $v$, and every further failure has probability $q$. Its expected remaining cost is $1/(1-q)$. On the whole $P$ failure range this continuation cost is the positive operator $H_P=(I_P-PQP)^{-1}$, and similarly $H_Q=(I_Q-QPQ)^{-1}$. These inverses exist because $c<1$. A first choice of $A$ with probability $r$ therefore has cost operator bounded below by
\begin{equation}\label{eq:detector-lower-operator}
 I+rPH_PP+(1-r)QH_QQ,
\end{equation}
with equality for subsequent alternation. On a principal block its largest eigenvalue is
\begin{equation}\label{eq:detector-block}
 1+\frac{1+\sqrt{1-4r(1-r)(1-q)}}{2(1-q)}.
\end{equation}
It increases with $q$ and is minimized uniquely at $r=1/2$. Unpaired parts have the corresponding zero-overlap bound, and the common kernel of $P,Q$ succeeds on the first call. The largest block thus gives \eqref{eq:detector-global}. If $c=0$, use a unit vector from each of the two nonzero orthogonal failure ranges; the same formula with $q=0$ applies.

With one live label, the same probability $r$ of choosing $A$ is used after every failure. A persistent shared seed would be extra memory and is not admitted. On a block, let $x,y$ be the remaining costs after $A,B$ fail. For $0<r<1$,
\[
 x=1+rx+(1-r)qy,\qquad y=1+rqx+(1-r)y.
\]
Solving and returning to the initial state gives $I+aP+bQ$ on that block, where
\[
 a=\frac{r/(1-r)+q}{1-q^2},\qquad
 b=\frac{(1-r)/r+q}{1-q^2}.
\]
The larger eigenvalue is $1+[a+b+\sqrt{q(a+b)^2+(1-q)(a-b)^2}]/2$. Since $a+b\ge2/(1-q)$, it is at least $1+(1+\sqrt q)/(1-q)=1+1/(1-\sqrt q)$, with equality at $r=1/2$. The deterministic endpoints have an undetectable failure state. The largest principal block proves \eqref{eq:detector-one}; the fair row realizes that operator bound in every block.

For the fixed-start unrestricted optimum, its first command row must be fair. On a largest principal block, both possible first failures have positive probability from the top vector of the fair cost operator. Every positive-probability repetition of the last failed detector strictly increases the expected cost there. Consequently an optimum needs a pure $B$ row after $A$ fails and a pure $A$ row after $B$ fails. These are two distinct rows, neither equal to the initial mixed row. Three live labels are necessary; those three rows and their obvious switching transitions suffice. This is a rigidity statement on the extremal block, not a claim about behavior on irrelevant physical subspaces.

The two-label parameter space is compact. For a sequence of proper two-label observers with uniformly bounded worst-state costs, every finite-prefix cost is continuous in the parameters and remains bounded at the limit. Monotone convergence gives a finite full cost there; unit costs imply proper operation. A minimizing sequence therefore has a proper minimizing limit. Such a limit cannot attain $V_*$ by the preceding row argument, proving the strict lower in \eqref{eq:detector-two}. Deterministic alternation from a fixed first detector uses two labels and has worst cost $1+(1-c^2)^{-1}$, proving the upper.

If the interface permits a randomized initial distribution, start fairly in the two pure switching labels. This stores the first choice in a counted label and attains $V_*$. One live label cannot attain it by \eqref{eq:detector-one}. The event/state counts follow from their exact halt semantics. Finally, a zero failure projection always succeeds. In finite dimension $\|PQ\|=1$ implies a nonzero vector in both failure ranges by the equality case for orthogonal projections. It remains there under every permitted command and is never detected. This proves the degenerate cases.
\end{proof}

\begin{corollary}[A primitive least-favorable preparation]\label{cor:detector-saddle}
Choose unit failure vectors $u,v$ in a largest principal-angle block, with $\langle u,v\rangle=c$, and put
\[
 \psi_*=(u+v)/\sqrt{2(1+c)},\qquad \rho_*=|\psi_*\rangle\langle\psi_*|.
\]
For $c=0$ choose any unit vectors in the two nonzero orthogonal failure ranges. This one pure state gives expected cost at least $V_*$ against every visible-history policy. Fair-first alternation has cost at most $V_*$ against every initial state. Thus it is an explicit minimax saddle without a policy-dependent preparation or a minimax exchange. The same preparation gives the one-live-label lower $V_1$ against every one-label row.
\end{corollary}
\begin{proof}
The first failure probability from $\rho_*$ is $(1+c)/2$ for either detector. Its post-failure state is respectively $u$ or $v$, whose least remaining cost is $1/(1-c^2)$ by the switching argument. Every first command mixture therefore costs at least $1+(1+c)/[2(1-c^2)]=V_*$. The uniform upper was proved above. For a one-label row, the cost at $\rho_*$ is $1+(a+b)(1+c)/2$. The bound $a+b\ge2/(1-c^2)$ gives $V_1$, attained by the fair row. This preparation is a witness for the stated detection task, not an automatically recurrent reference row after arbitrary rejected experiments.
\end{proof}

\subsection{The exact two-live-label frontier}
Put $q=c^2$. For $0<r\le1$ define
\begin{align}
 a_q(r)&=r+q,&b_q(r)&=(1-r)(r^{-1}+q),\label{eq:two-label-ab}\\
 F_q(r)&=1+\frac{a_q(r)+b_q(r)+\sqrt{(a_q(r)-b_q(r))^2+4q a_q(r)b_q(r)}}{2(1-q^2)}.
 \label{eq:two-label-F}
\end{align}

\begin{theorem}[Exact two-label optimum and a golden-ratio transition]\label{thm:two-label-exact}
In the nondegenerate class above,
\begin{equation}\label{eq:two-label-opt}
 V_2(c)=\min_{0<r\le1}F_{c^2}(r).
\end{equation}
A minimizing observer has a fixed initial label $S$ choosing $A$ with probability $r$ and $B$ otherwise, and a label $T$ choosing $B$ purely. After an $A$ failure it uses $T$; after a $B$ failure it uses $S$. Success commits immediately. It realizes the optimum for every finite-dimensional pair with the same largest overlap $c$.

Let $q_*=(\sqrt5-1)/2$. The minimizer is unique and is characterized by
\begin{equation}\label{eq:two-label-transition}
 \begin{array}{c|c|c}
 q&r_{\rm opt}&V_2\\\hline
 0&q_*&(1+\sqrt5)/2\\
 0<q<q_*&\text{unique }r\in(0,1)\text{ with }F'_q(r)=0&F_q(r_{\rm opt})\\
 q_*\le q<1&1&1+(1-q)^{-1}.
 \end{array}
\end{equation}
In the interior case, put $a=a_q(r)$, $b=b_q(r)$, $D=\sqrt{(a-b)^2+4qab}$ and $b'=-r^{-2}-q$. The unique equation is explicitly
\begin{equation}\label{eq:two-label-stationary}
 (1+b')D+(a-b)(1-b')+2q(b+ab')=0.
\end{equation}
Squaring does not license extraneous roots. Together with the one-label and unrestricted laws, this determines the exact fixed-start cost at every integer live-memory budget: $V_1$ at one, $V_2$ at two, and $V_*$ at three or more. State-stopped total counts add one terminal label; event-stopped counts do not.
\end{theorem}
\begin{proof}
First prove a converse against every two-label controller, including randomized post-failure transitions. Restrict the initial state to a largest principal block with failure states $u,v$ and squared overlap $q$. For $q=0$ use any unit vectors in the two nonzero orthogonal ranges. Relabel detector names if necessary so that the initial label has the larger probability $r$ of choosing $A$. The other label has probability $t$, with $0\le t\le r\le1$. Proper worst-state operation requires $r>0$ and $t<1$.

Even allowing the controller to choose whichever of its two labels it prefers after a failure, the least remaining costs at $u,v$ are
\begin{equation}\label{eq:two-row-values}
 x=\frac{(1-t)^{-1}+q/r}{1-q^2},\qquad
 y=\frac{r^{-1}+q/(1-t)}{1-q^2}.
\end{equation}
Here is a direct verification for arbitrary subsequent choices, not an assumed stationary optimum. One has $x-qy=(1-t)^{-1}$ and $y-qx=r^{-1}$. For any command probability $z\in[t,r]$,
\[
 1+zx+(1-z)qy-x=\frac{z-t}{1-t}\ge0,
 \qquad
 1+zqx+(1-z)y-y=1-\frac zr\ge0.
\]
These are lower potential inequalities for every allowed label and randomized label choice. Iteration to a finite horizon bounds actual expected calls below by the initial potential minus a remaining term bounded by $\max(x,y)$ times survival. The remainder vanishes for a proper policy. An improper policy has infinite call cost. Thus $x,y$ are universal lower bounds. They are attained by using the $t$-row after an $A$ failure and the $r$-row after a $B$ failure, as their equations verify.

From the prescribed initial label, the first call therefore has cost operator at least $I+rxP+(1-r)yQ$ on this block. For fixed $r$, both $x,y$ are nondecreasing in $t$, so this operator is minimized at $t=0$. Its coefficients then are
\[
 rx=\frac{r+q}{1-q^2},\qquad
 (1-r)y=\frac{(1-r)(r^{-1}+q)}{1-q^2}.
\]
Its largest eigenvalue is $F_q(r)$. This proves the lower without assuming deterministic transitions, a prior, or an uncounted initial seed.

For the matching upper use the displayed $S,T$ observer. Its exact block costs are \eqref{eq:two-row-values} at $t=0$. In a smaller principal block with squared overlap $q'\le q$, both nonnegative coefficients $a_{q'}(r)/(1-q'^2)$ and $b_{q'}(r)/(1-q'^2)$ increase with $q'$. The larger eigenvalue of $aP+bQ$ is nondecreasing in either nonnegative coefficient by positive matrix order, and in overlap by its square-root formula. Thus the worst block is the largest. Unpaired ranges have no greater than the corresponding zero-overlap bound, and the common kernel succeeds immediately. The same observer therefore has cost at most $F_q(r)$ for every initial state and attains it on the largest block. Its finite post-failure costs prove proper operation. Equality in \eqref{eq:two-label-opt} follows.

For $q=0$, $F_0(r)=1+\max\{r,r^{-1}-1\}$. The increasing and decreasing terms meet uniquely at $r^2+r-1=0$, giving the first line of \eqref{eq:two-label-transition}.

For $0<q<1$, use one fixed principal block. Its numerator matrix is $a_q(r)P+b_q(r)Q$, with $a_q$ affine and $b_q''(r)=2/r^3>0$. Its largest eigenvalue is strictly convex. Indeed the convex-combination matrix at two different parameters exceeds the matrix at the combined parameter by a positive multiple of $Q$. A top eigenvector of the latter has positive $Q$-expectation since $q>0$ and $a_q>0$. Evaluation at that vector proves strict convexity. The top eigenvalue is simple, so the function is differentiable, and it diverges as $r\downarrow0$.

At $r=1$ the numerator is $(1+q)P$. Its top vector is $u$, with $Q$-expectation $q$. Differentiation at this simple eigenvector gives
\[
 (1-q^2)F'_q(1)=1-q-q^2.
\]
For $q\ge q_*$ this is nonpositive, so strict convexity makes $r=1$ the unique minimum. For $q<q_*$ it is positive; divergence at zero and strict convexity give a unique interior stationary point. Differentiating the square root gives \eqref{eq:two-label-stationary}. This proves every case. The already proved live-role counts give the full integer frontier.
\end{proof}

The golden-ratio threshold concerns the squared principal-angle cosine and the optimal controller row, not a universal physical constant. At small overlap one continuing row is pure and the initial row is mixed. At and above the threshold the optimal two-label policy is fixed-start deterministic alternation. A third live label restores a fair initial choice without contaminating either continuing row and strictly improves the worst-state cost.

\subsection{A sharp memory-constrained minimax gap}
The exact finite-memory frontier also identifies a limitation of replacing a uniform converse by an optimized single-prior calculation. Let $\mathcal C_2$ be the proper fixed-start observers with at most two live labels in this detection interface.

\begin{corollary}[Minimax gap at two live labels]\label{cor:memory-minimax-gap}
For $P,Q\ne0$ and $c<1$,
\begin{equation}\label{eq:memory-minimax-gap}
 \inf_{\pi\in\mathcal C_2}\sup_\rho\E_\rho^\pi\tau
       =V_2(c)>V_*(c)
       =\sup_\rho\inf_{\pi\in\mathcal C_2}\E_\rho^\pi\tau.
\end{equation}
In particular no single initial preparation supplies the sharp two-label robust lower $V_2$. At zero overlap the gap is $(\sqrt5-2)/2$. Permitting a randomized initial distribution over the two counted pure labels removes this gap and gives robust value $V_*$. No analogous gap occurs at one live label or at three or more for this experiment.
\end{corollary}
\begin{proof}
The left equality and strict inequality are the memory theorems. For any fixed initial state, choose the better of the two deterministic first detectors and then alternate. Both choices are proper two-label observers. Their better cost is no larger than their average, which is the cost of fair-first alternation and is at most $V_*$. Hence the maximin is at most $V_*$. The single state $\rho_*$ in Corollary~\ref{cor:detector-saddle} gives cost at least $V_*$ against every policy, proving the reverse inequality. When $c=0$, $V_2=(1+\sqrt5)/2$ and $V_*=3/2$, giving the displayed gap. A randomized initial label realizes fair-first alternation directly. For one label, the same least-favorable state and fair independent row form a saddle at $V_1$; for three or more the full-history saddle applies.
\end{proof}

The policy in the inner maximin may depend on the fixed preparation; the robust policy on the left may not. The fixed-start two-label class is not closed under the convexification that would mix the two pure initial rows while remembering which continuing row to use. A free random initial distribution changes precisely that resource. Thus a minimax exchange would erase a real memory constraint. This is an exact result for the stated detector model, not a failure of the general sufficient transfer theorem or a claim that every finite-memory problem has such a gap.
